\section{Search in a 2D Matrix II}
\label{sec:bs-search-2d-matrix-2}

\problemstatement{We are given an $n \times m$ matrix in which every row
is sorted in increasing order \textbf{left to right}, and every column is
sorted in increasing order \textbf{top to bottom} -- but, unlike
Section~\ref{sec:bs-search-2d-matrix}, there is \textbf{no} guarantee
chaining one row to the next (the last element of a row may be larger
than the first element of the next row). Given a target $X$, determine
whether it exists in the matrix.}

\begin{patternbox}
The missing row-chaining guarantee is the keyword to notice by its
\emph{absence}: it means the matrix can \textbf{not} be treated as a
flattened sorted 1D array (Section~\ref{sec:bs-search-2d-matrix}'s
virtual-index trick would give wrong answers here, since the virtual
sequence is not globally sorted). Instead, the two \emph{independent}
sortedness guarantees -- rows increasing rightward, columns increasing
downward -- call for a different technique: starting from a corner where
one direction increases and the other decreases the candidate's relation
to $X$, and walking a single staircase path.
\end{patternbox}

\subsection*{Understanding the problem carefully}

Start at the \textbf{top-right} corner, $(\textit{row}, \textit{col}) =
(0, m-1)$. From this corner: moving \textbf{left} (decreasing
\textit{col}) only ever decreases the value (row sortedness), and moving
\textbf{down} (increasing \textit{row}) only ever increases the value
(column sortedness). This gives us a clean three-way comparison at every
step:

\begin{itemize}
  \item If $\textit{matrix}[\textit{row}][\textit{col}] = X$: found it.
  \item If $\textit{matrix}[\textit{row}][\textit{col}] > X$: every entry
        further right in this row is even larger (row sortedness), so $X$
        cannot be in the rest of this row; move left:
        $\textit{col} \gets \textit{col}-1$.
  \item If $\textit{matrix}[\textit{row}][\textit{col}] < X$: every entry
        further up in this column is even smaller (column sortedness), so
        $X$ cannot be above in this column; move down:
        $\textit{row} \gets \textit{row}+1$.
\end{itemize}

\begin{ideabox}[Why the Top-Right Corner Is the Only Safe Starting Point]
The top-left corner is the smallest element in the entire matrix (both
row and column sortedness make every other element $\ge$ it), so both
``move right'' and ``move down'' would increase the value -- a comparison
there gives no way to discard a direction safely. The bottom-right corner
is symmetric (the largest element, both directions only decrease).
\emph{Only} the top-right (or, symmetrically, bottom-left) corner has the
property that one direction strictly increases the value while the other
strictly decreases it, which is exactly what is needed to safely discard
an entire row or column on each comparison.
\end{ideabox}

\subsection*{Why this staircase walk is correct and terminates in $O(n+m)$}

At each step, the walk either eliminates an entire remaining row's worth
of values (when it moves left, the eliminated value and everything to its
right in that row are confirmed too large) or an entire remaining column's
worth (when it moves down, the eliminated value and everything above it in
that column are confirmed too small) -- but more precisely, each step
moves the pointer one position left or one position down, and since
\textit{col} can decrease at most $m$ times and \textit{row} can increase
at most $n$ times before going out of bounds, the walk terminates after at
most $n+m$ steps, either finding $X$ or running off the matrix entirely
(confirming $X$ is absent).

\subsection*{Algorithm}

\begin{enumerate}
  \item Initialize $\textit{row} \gets 0$, $\textit{col} \gets m-1$.
  \item While $\textit{row} < n$ and $\textit{col} \ge 0$:
  \begin{enumerate}
    \item If $\textit{matrix}[\textit{row}][\textit{col}] = X$: return
          \texttt{true}.
    \item If $\textit{matrix}[\textit{row}][\textit{col}] > X$:
          $\textit{col} \gets \textit{col}-1$.
    \item Else: $\textit{row} \gets \textit{row}+1$.
  \end{enumerate}
  \item Return \texttt{false}.
\end{enumerate}

\begin{algorithm}[H]
\caption{Search in a 2D Matrix II (Staircase Search)}
\begin{algorithmic}[1]
\Procedure{SearchMatrixII}{\textit{matrix}, $X$}
  \State $\textit{row} \gets 0$, $\textit{col} \gets m-1$
  \While{$\textit{row} < n$ \textbf{and} $\textit{col} \ge 0$}
    \If{$\textit{matrix}[\textit{row}][\textit{col}] = X$} \Return \textbf{true} \EndIf
    \If{$\textit{matrix}[\textit{row}][\textit{col}] > X$}
      \State $\textit{col} \gets \textit{col} - 1$
    \Else
      \State $\textit{row} \gets \textit{row} + 1$
    \EndIf
  \EndWhile
  \State \Return \textbf{false}
\EndProcedure
\end{algorithmic}
\end{algorithm}

\subsection*{Dry run}

\dryrun{Take $\textit{matrix} = \begin{bmatrix} 1&4&7&11 \\ 2&5&8&12 \\
3&6&9&16 \\ 10&13&14&17 \end{bmatrix}$, $X = 5$. ($n=m=4$.)}

\begin{itemize}
  \item $\textit{row}=0,\textit{col}=3$: $\textit{matrix}[0][3]=11 > 5$.
        $\textit{col}=2$.
  \item $\textit{row}=0,\textit{col}=2$: $\textit{matrix}[0][2]=7 > 5$.
        $\textit{col}=1$.
  \item $\textit{row}=0,\textit{col}=1$: $\textit{matrix}[0][1]=4 < 5$.
        $\textit{row}=1$.
  \item $\textit{row}=1,\textit{col}=1$: $\textit{matrix}[1][1]=5 = 5$.
        Return \texttt{true}.
\end{itemize}

\subsection*{C++ implementation}

\begin{lstlisting}
#include <bits/stdc++.h>
using namespace std;

bool searchMatrixII(vector<vector<int>>& matrix, int X) {
    int n = matrix.size(), m = matrix[0].size();
    int row = 0, col = m - 1;

    while (row < n && col >= 0) {
        if (matrix[row][col] == X) return true;
        else if (matrix[row][col] > X) col--;
        else row++;
    }
    return false;
}

int main() {
    vector<vector<int>> matrix = {
        {1, 4, 7, 11},
        {2, 5, 8, 12},
        {3, 6, 9, 16},
        {10, 13, 14, 17}
    };
    cout << (searchMatrixII(matrix, 5) ? "true" : "false") << endl;
    return 0;
}
\end{lstlisting}

\begin{complexitybox}
\textbf{Time Complexity.} Each step decreases \textit{col} or increases
\textit{row}, and there are at most $m$ decreases and $n$ increases before
the pointer exits the matrix, giving
\[
  O(n+m).
\]
Note this is not $O(\log(nm))$ -- the lack of row-chaining genuinely costs
us the logarithmic complexity available in
Section~\ref{sec:bs-search-2d-matrix}, even though this is still a vast
improvement over the $O(nm)$ brute-force scan.
\textbf{Space Complexity.} $O(1)$ auxiliary space.
\end{complexitybox}

\begin{takeawaybox}
When a 2D matrix is sorted along both rows and columns independently but
not chained into one global order, look for a corner where the two
sortedness directions act as \emph{opposite} comparisons (one direction
increasing, the other decreasing) relative to the target -- that corner is
always the unique safe starting point for a linear ``staircase'' walk,
even though it does not give the $O(\log(nm))$ that full chaining would
allow.
\end{takeawaybox}
